\section{Recommended Reading}

\begin{readingbox}
\textbf{Foundation}
\begin{itemize}[nosep]
  \item Eugene Izhikevich, \emph{Dynamical Systems in Neuroscience}, for a
        rigorous but usable route from nonlinear state dynamics to intervention
        thinking.
  \item Steven Strogatz, \emph{Nonlinear Dynamics and Chaos}, for intuition
        about regime change, local stability, and why small changes near a
        boundary can matter.
  \item John Holland, \emph{Emergence}, for the idea that repeated local rules
        can accumulate into durable higher-level structure.
\end{itemize}
\textbf{Modern Research}
\begin{itemize}[nosep]
  \item \emph{Dynamical structure-function correlations provide robust and
        generalizable signatures of consciousness in humans} (Communications
        Biology, 2024), as support for treating conscious states as structured
        dynamical regimes rather than undifferentiated ``motivation.''
  \item \emph{Cortical connectivity, local dynamics and stability correlates of
        global conscious states} (Communications Biology, 2025), for the link
        between stability structure and global state organization.
  \item \emph{Falling asleep follows a predictable bifurcation dynamic}
        (Nature Neuroscience, 2025), for a concrete example of state transition
        analysis that strengthens window-sensitive intervention thinking.
  \item \emph{Metastability demystified: the foundational past, the pragmatic
        present and the promising future} (Nature Reviews Neuroscience, 2024),
        for the medium-lane vocabulary behind temporary stabilization and
        flexible switching.
  \item \emph{Associative conditioning in gene regulatory network models
        increases integrative causal emergence} (Communications Biology, 2025),
        as a bridge source for history dependence and emergence-informed
        intervention design.
\end{itemize}
\textbf{Frontier}
\begin{itemize}[nosep]
  \item Frontier consciousness-physics proposals should be read only as
        hypothesis-generating extensions. They may inspire alternative system
        pictures, but they do not currently justify operational intervention
        rules in this chapter.
\end{itemize}
\end{readingbox}

\begin{summarybox}
This chapter turns the book's research spine into operational design. The main
practical lesson is not ``try harder.'' It is: diagnose the transition, choose
the cheapest controllable variable, exploit a real decision window, and prefer
designs that create better future history rather than one-off acts of effort.
The strong lane justifies transition-sensitive thinking; the medium lane helps
explain why repetition and local structure matter; the final protocol remains a
design inference whose value must be checked in use.
\end{summarybox}
